\scanpage{10}
\mspage{8}
\uline{Remarques} 1) L'ens des axiomes $F_1:F_5$ est trivialement satisfait si $S=\varnothing$ (i.e. si $L_r=\varnothing$ i.e. tous les lieux sont singuliers).
Alors la structure envisagée se réduit à celle d'un ens $\unsure{L}$.
On sera amené tantôt \unsure{d'enrichir|à enrichir} un peu la structure, en introduisant (ds le cas général) une relation d'équivalence sur $L_s$\ldots

2) Soit $L'\subset L$, partie quelc. de $L$, et soit $\unsure{S'|S^1}=\{I\in S\mid\tilde I\subset L'\}$. On a donc
\[
S'\longrightarrow\mathcal P(L'),\qquad I\unsure{\mapstochar\relbar|\mapsto}\tilde I
\]
\[
S'\longrightarrow\mathcal P_2(L'),\qquad I-\delta I.
\]
Ceci dit, pour cette structure, $(\unsure{L'},S')$ est un modèle de forme de dim 1, i.e. les axiomes $F_1-F_5$ sont satisfaits.

Cas particulier: soit $I\in S$, $L'=\tilde I$, alors $S'$ est l'ens \inserted{$S_I$} des segments tels que $\unsure{\tilde I|\tilde I'}\subset I$, et s'identifie donc à une partie de $\mathcal P_2(L')$.
\mspage{9}

\rotatebox{90}{\reflectbox{\unsure{EXTRA STRONG}}}
\rotatebox{90}{\reflectbox{\unsure{EXTRA STRONG}}}
